\section{Introdução}

O ônibus é o principal meio de transporte coletivo na maior parte das cidades brasileiras, e a
confiabilidade do serviço pesa diretamente na decisão de usá-lo ou não. Quando o passageiro não
consegue prever quanto tempo vai esperar ou quanto vai durar a viagem, ele passa a reservar
mais tempo para cada deslocamento ou migra para outros modos, o que vai na contramão das
políticas de mobilidade de baixo carbono discutidas na literatura \cite{WIMBADI2021103023}.
A percepção de qualidade do transporte público depende, entre outros fatores, da pontualidade e
da regularidade do serviço \cite{https://doi.org/10.1002/atr.106}.

A disponibilidade de dados abertos de geolocalização dos veículos tornou possível medir essa
previsibilidade com frequência e em larga escala. \cite{courb} usaram
dados de GPS dos ônibus de Curitiba e Modelos Mistos para estudar o tempo de deslocamento entre
paradas consecutivas de uma linha, mostrando que a maior parte da variabilidade está associada
ao trecho percorrido e que há dependência temporal entre observações sucessivas. Na mesma
direção, \cite{10.1109/TITS.2022.3145243} tratam a
previsibilidade das operações de transporte como algo que pode ser modelado e acompanhado, e
não apenas como pontualidade.

Este trabalho replica a metodologia de \cite{courb} para Belo Horizonte,
usando como caso de estudo a linha 4110, no sentido Belvedere $\rightarrow$ Dom Cabral. Além
da réplica, o trabalho propõe uma extensão: comparar os horários de passagem observados com o
quadro de horários publicado no GTFS da BHTRANS, para medir o atraso em relação ao programado e
verificar se ele varia com a faixa horária, o tipo de dia e a posição da parada no percurso.

As contribuições do trabalho são: (i) a aplicação da metodologia de Curitiba a outra capital,
com outra fonte de dados, o que permite avaliar a transferibilidade do método; (ii) a
validação fora da amostra dos modelos, com métricas quantitativas; (iii) a medição da
diferença entre o tempo de viagem real e o programado na linha estudada; e (iv) a
disponibilização do código e dos dados utilizados, permitindo a reprodução dos resultados.

O restante do texto está organizado da seguinte forma. A Seção~\ref{sec:relacionados}
apresenta os trabalhos relacionados. A Seção~\ref{sec:metodologia} descreve a coleta, o
tratamento dos dados e os modelos. A Seção~\ref{sec:resultados} apresenta e discute os
resultados, e a Seção~\ref{sec:conclusao} traz as conclusões, as limitações e os trabalhos
futuros.
